% ============================================================
% DRAFT ONLY — not meant to compile, nothing \input's this file.
% This supersedes introduction_new_GAs_addition.tex — the GA
% paragraph is folded into the tightened Background below, so just
% use this file now and ignore that one.
%
% WHAT THIS IS: full replacement text for \mysection{Background} and
% \mysection{Scope and Limitations} in introduction.tex, tightened per
% your request. Lines starting with % are my notes/the CURRENT text
% for reference — DO NOT paste those. Lines with no % prefix are the
% REPLACE-WITH text — paste those in directly, paragraph breaks as
% shown.
%
% What changed and why, so nothing feels like it vanished on you:
% - Background: cut restatement (the "every command sent... every
%   camera frame..." privacy point was made twice back to back — now
%   once), merged the Jetson rejection from four numbered "reasons"
%   down to three real ones stated once each instead of twice, and
%   folded in the new-GA paragraph (GA5/GA9) we drafted earlier so you
%   only have one file to paste from. Nothing was deleted that wasn't
%   a repeat of something still there.
% - Scope and Limitations: same approach — merged two of the six
%   paragraphs that were both just "X is out of scope" with their own
%   throat-clearing into one paragraph, and fixed a few grammar slips
%   ("does not perform no part", "Rather than guaranteeing" as a
%   stray sentence fragment) while I was in there.
% - Section headers and \label{} lines are UNCHANGED — only the body
%   paragraphs underneath them are replaced.
% ============================================================


% ---------------- BACKGROUND ----------------
% CURRENT: everything under \mysection{Background}{Background} /
% \label{sec:background}, i.e. all paragraphs up to (not including)
% \mysection{Problem Statement and Research Questions}.
%
% REPLACE WITH (paste as the full body of the Background section):

A robot that takes spoken instructions removes two obstacles that make most mobile robots
difficult to operate day to day: a screen to navigate on, and an operator trained on its
command set. This matters even more for a small indoor rover, since tasks such as checking
whether a door is open or looking at what is on a shelf are exactly the kind of thing a person
would rather say out loud than type into an application.

\Acp{llm} are transformer-based neural networks trained on large collections of text to predict
the next word given what came before, which makes them well suited to interpreting everyday
instructions and gives them a broad general-knowledge base. Given a natural-language command
and a description of what a robot can do, an \ac{llm} can break the command into a sequence of
simple actions, a \ac{robotactionprimitive} (RAP), built only from actions the robot can
actually perform, without a developer hand-coding a fixed list of recognised
phrases~\cite{ahn2022saycan}. Ambiguous commands can be resolved the same way, through
clarifying questions: when a command leaves out information, the model asks the user
follow-up questions before committing to a plan. This back-and-forth is referred to in this
report as a \emph{multi-turn dialogue}~\cite{hori2025longhorizon}.

Handing planning to a language model introduces two problems. Spoken commands are often
ambiguous, and a system that always acts on its first interpretation will act on the wrong one
often enough to matter~\cite{ren2023knowno}. Language models also hallucinate: they produce
fluent, confident output that does not correspond to anything real, and this becomes more
likely as a task is broken into more steps~\cite{liang2024introspective}. The techniques that
catches ambiguity or hallucination before a plan is set can also run while the plan is carried
out, by monitoring the robot's observations for signs that something has gone
wrong~\cite{sinha2024aesop, elhafsi2023semantic}.

Systems that demonstrate these techniques typically rely on large cloud-hosted models and
robots with rich sensing such as \acs{lidar}, an \ac{imu} and wheel encoders. Both assumptions
are reasonable in a research laboratory, but do not hold in every setting where a
voice-commanded robot would be useful.

A private home is one such setting. Every command sent to a cloud-hosted model, and every
camera frame used to check progress, leaves the home for a third party's server — unacceptable
to many households for a device that moves through their living space and records it. A home
robot also has to be affordable, which rules out \acs{lidar} and a powerful on-board computer,
leaving a camera as the most practical sensor. Low-cost drive platforms also typically lack
wheel encoders and an \ac{imu}, and slip enough that a command such as `drive forward one
metre' cannot be trusted to be carried out accurately. In this setting the language models
must run on hardware the household already owns, such as a desktop with a consumer graphics
card, and the robot must rely on its camera to know whether its plan is working. Where privacy
is not a concern and a reliable connection is available, a cloud-hosted model remains the
better choice, since it is larger and more capable than anything that fits on consumer
hardware.

An alternative would be to mount a dedicated embedded AI accelerator, such as an NVIDIA Jetson
Orin Nano or Orin NX, directly on the rover and run the whole pipeline on board. This project
deliberately does not do that. First, the home-robot framing above assumes the household
already owns the compute the system needs, in the form of a desktop or laptop. This requires a
separate accelerator module reintroduces exactly the extra, specialised hardware cost this
project is trying to avoid. Second, keeping inference off the rover, on a host machine that can
run the same planner locally or swap to a cloud API, keeps the local-versus-cloud comparison
closer to apples-to-apples on model capability alone. If the rover itself hosted the model,
switching the benchmark to a cloud API would also change the hardware story. Third, it keeps
the result general: the rover's own on-board computer, a Raspberry Pi, only has to stream
camera frames and relay drive and sensor commands over the local network, so the same rover
would work unchanged with a different host machine entirely.

This framing also lines up with Graduate Attribute 5 (Use of Engineering Tools), which shows the
limitations of the modern engineering and IT tools applied to a problem, not just how to
operate them. Exactly what RQ1 and RQ2 set out to do for a locally hosted LLM and VLM
pipeline. Building that pipeline from scratch, with no existing local reference implementation
to follow, also draws on Graduate Attribute 9 (Independent Learning Ability), in its new
emphasis on adapting to emerging technologies.


% ---------------- SCOPE AND LIMITATIONS ----------------
% CURRENT: everything under \mysection{Scope and Limitations}{Scope
% and Limitations} / \label{sec:scope-and-limitations}, i.e. all
% paragraphs up to (not including) \mysection{Report Outline}.
%
% REPLACE WITH (paste as the full body of Scope and Limitations):

This project is concerned with turning vague spoken commands into robot action plans, so its
scope is limited to the pipeline from speech to executed plan: transcription, clarification of
ambiguous commands through spoken dialogue, plan generation and verification, and the spoken
confirmation step. Navigation is used only to carry out and test the plans; improving
navigation itself is not an aim of this project.

Before the rover is used in a room, the room is recorded once in a short walkthrough video,
which a \ac{vlm} describes frame by frame to produce a written summary of the objects present
and their layout. This summary is given to the planner, so plans refer only to things that are
actually present, and selected frames are stored alongside their descriptions as a reference
set that the rover's live camera frames are compared against during a mission. The scan is not
turned into a three-dimensional model or map; building a general-purpose mapping system is
outside this project's scope.

Cloud-hosted models are used only as a benchmark for the locally hosted pipeline and are not
part of the system under test. The rover itself is deliberately limited to a camera, an
ultrasonic distance sensor and a line sensor, with no wheel encoders, \ac{imu} or \acs{lidar}. This
reflects the low-cost home robot described in Section~\ref{sec:background} and forces the
pipeline to rely on the camera to judge whether a plan is succeeding. As a result, movement is carried out without precise measurement of
distance or heading, and commands requiring distances beyond what the ultrasonic sensor can
measure are not supported.

In the same spirit, the rover's own on-board computer is treated as instrumentation rather
than part of the system under test: it streams camera frames to the host machine and relays
back the drive and sensor commands it receives, performing none of the planning, perception or
clarification itself. The language and vision models that answer RQ1 and RQ2 run on the host
machine defined in Section~\ref{sec:background}, so both questions probe those models'
capability on fixed, known hardware rather than being confounded by whatever embedded computer
happens to be fitted to the rover.

Similarly, the system is not designed to operate in arbitrary rooms or with arbitrary commands:
evaluation is limited to a small number of test rooms and a fixed set of commands at each level
in Table~\ref{tab:command-levels}, sufficient to answer RQ1 and RQ2. The results describe the
pipeline's behaviour on the cases tested, not a guarantee of its performance in other rooms or
with other commands.
